\documentclass[11pt,a4paper]{article}

\usepackage[utf8]{inputenc}
\usepackage[T2A]{fontenc}
\usepackage[ukrainian,english]{babel}
\usepackage[top=22mm,bottom=22mm,left=22mm,right=22mm]{geometry}
\usepackage{hyperref}
\usepackage{fancyhdr}
\usepackage{parskip}
\usepackage{tcolorbox}
\usepackage{enumitem}
\usepackage{listings}
\usepackage{xcolor}

\tcbuselibrary{skins,breakable}

\hypersetup{
    pdftitle={Workshop 07 — MCP Servers: Handout},
    pdfauthor={Vadym Bondarenko},
    colorlinks=true,
    linkcolor=blue,
    urlcolor=cyan,
}

\pagestyle{fancy}
\fancyhf{}
\fancyhead[L]{\small\textit{Workshop 07 — MCP Servers}}
\fancyhead[R]{\small\thepage}
\fancyfoot[C]{\small\textit{vcdev.me \textperiodcentered{} 2026}}
\renewcommand{\headrulewidth}{0.4pt}

\lstdefinelanguage{json}{
    basicstyle=\ttfamily\small,
    sensitive=true,
    morestring=[b]",
    comment=[l]{//},
}

\lstdefinelanguage{yaml}{
    keywords={true,false,null},
    sensitive=false,
    comment=[l]{\#},
    morestring=[b]',
    morestring=[b]"
}

\lstset{
    basicstyle=\ttfamily\small,
    breaklines=true,
    frame=single,
    backgroundcolor=\color{gray!10},
    rulecolor=\color{gray!40},
    xleftmargin=8pt,
    xrightmargin=8pt,
    aboveskip=6pt,
    belowskip=6pt,
}

\title{%
  {\Huge\bfseries Workshop 07: MCP-сервери}\\[0.6em]
  {\large Будуємо свій на Python --- handout}
}
\author{Vadym Bondarenko \\ \small\href{mailto:vadym.bondarenko@vcdev.me}{vadym.bondarenko@vcdev.me}}
\date{2026}

\begin{document}

\maketitle

\begin{abstract}
Пост-воркшоп референс для самостійного перечитування. Не повторює слайди --- сухе зведення команд, API, чек-листів. Усі факти посилаються на офіційні docs (\href{https://modelcontextprotocol.io}{modelcontextprotocol.io}, \href{https://code.claude.com/docs/en/mcp}{code.claude.com}).
\end{abstract}

\tableofcontents
\newpage

\section{Що таке MCP}

Model Context Protocol --- open-source протокол між AI-host-додатком і твоїми тулами/даними. Client--server, JSON-RPC 2.0.

\subsection{Учасники}

\begin{tabular}{ll}
\hline
\textbf{Роль} & \textbf{Що} \\
\hline
MCP Host & AI-додаток (Claude Code, Claude Desktop, VS Code) \\
MCP Client & Компонент усередині host, по одному на сервер \\
MCP Server & Програма, яка експонує контекст / тули \\
\hline
\end{tabular}

\textbf{Local server:} STDIO транспорт, subprocess host-а. \textbf{Remote server:} Streamable HTTP, окрема машина.

\subsection{Два рівні}

\begin{itemize}
\item \textbf{Data layer} --- JSON-RPC 2.0; lifecycle, primitives, notifications. Однаковий для всіх транспортів.
\item \textbf{Transport layer} --- STDIO або Streamable HTTP. Framing, з'єднання, auth.
\end{itemize}

\section{Примітиви сервера}

\begin{tabular}{lll}
\hline
\textbf{Примітив} & \textbf{Хто викликає} & \textbf{Методи} \\
\hline
\texttt{tools} & Модель (з апруву) & \texttt{tools/list}, \texttt{tools/call} \\
\texttt{resources} & Юзер/host & \texttt{resources/list}, \texttt{resources/read} \\
\texttt{prompts} & Юзер через slash & \texttt{prompts/list}, \texttt{prompts/get} \\
\hline
\end{tabular}

\textbf{Notifications:} сервер може push-ити \texttt{notifications/tools/list\_changed} тощо --- client рефрешить state без запиту.

\section{Python SDK: setup}

\textbf{System requirements:} Python 3.10+, MCP SDK 1.2.0+.

\begin{lstlisting}[language=bash]
# Встановити uv
curl -LsSf https://astral.sh/uv/install.sh | sh

# Створити проєкт
uv init my-server
cd my-server
uv venv
source .venv/bin/activate
uv add "mcp[cli]"
\end{lstlisting}

\section{FastMCP: декоратори}

\subsection{Tool}

\begin{lstlisting}[language=python]
from mcp.server.fastmcp import FastMCP

mcp = FastMCP("name")

@mcp.tool()
def add(a: int, b: int) -> int:
    """Add two numbers."""
    return a + b
\end{lstlisting}

\textbf{Type hints} $\to$ JSON Schema \texttt{inputSchema}. \textbf{Docstring} $\to$ description. Можна \texttt{async def}.

\subsection{Resource}

\begin{lstlisting}[language=python]
@mcp.resource("file://notes")           # static
def get_notes() -> str:
    return open("notes.md").read()

@mcp.resource("file://docs/{name}")     # templated
def get_doc(name: str) -> str:
    return open(f"docs/{name}.md").read()
\end{lstlisting}

\textbf{Static} resources показуються у \texttt{resources/list}. \textbf{Templated} --- ні (бо нескінченно багато), але client може запросити конкретний URI.

\subsection{Prompt}

\begin{lstlisting}[language=python]
@mcp.prompt(title="Code Review")
def review_code(code: str) -> str:
    return f"Please review this code:\n\n{code}"

# З роlями:
from mcp.server.fastmcp.prompts import base

@mcp.prompt()
def debug_session(error: str) -> list[base.Message]:
    return [
        base.UserMessage("I'm seeing this error:"),
        base.UserMessage(error),
    ]
\end{lstlisting}

\section{Запуск і транспорти}

\subsection{STDIO (default)}

\begin{lstlisting}[language=python]
if __name__ == "__main__":
    mcp.run(transport="stdio")
\end{lstlisting}

\textbf{Критично:} жодного \texttt{print()} без \texttt{file=sys.stderr}. Stdout зарезервований для JSON-RPC framing.

\begin{lstlisting}[language=python]
import sys, logging
logging.basicConfig(level=logging.INFO, stream=sys.stderr)
print("debug", file=sys.stderr)   # OK
\end{lstlisting}

\subsection{Streamable HTTP}

\begin{lstlisting}[language=python]
mcp.run(transport="streamable-http")
# default: http://localhost:8000/mcp
\end{lstlisting}

\textbf{SSE (deprecated)} --- є для legacy серверів, але не для нових проєктів.

\section{Підключення до Claude Code}

\subsection{CLI: \texttt{claude mcp add}}

\textbf{Правило ordering:} опції (\texttt{--transport}, \texttt{--env}, \texttt{--scope}, \texttt{--header}) йдуть \textbf{до} імені; \texttt{--} відокремлює name від command.

\begin{lstlisting}[language=bash]
# Local stdio
claude mcp add --transport stdio --env API_KEY=xxx myserver \
  -- uv run server.py

# Remote HTTP (Bearer)
claude mcp add --transport http myserver https://api.example.com/mcp \
  --header "Authorization: Bearer your-token"

# Remote HTTP (OAuth)
claude mcp add --transport http myserver https://mcp.notion.com/mcp
# далі усередині сесії: /mcp -> auth flow
\end{lstlisting}

\subsection{\texttt{.mcp.json} (project scope)}

\begin{lstlisting}[language=json]
{
  "mcpServers": {
    "my-clock": {
      "command": "uv",
      "args": ["--directory", "/abs/path", "run", "server.py"],
      "env": {"DEBUG": "1"}
    },
    "remote-api": {
      "type": "http",
      "url": "${API_BASE_URL:-https://api.example.com}/mcp",
      "headers": {
        "Authorization": "Bearer ${API_KEY}"
      }
    }
  }
}
\end{lstlisting}

Підтримує \texttt{\$\{VAR\}} і \texttt{\$\{VAR:-default\}} у \texttt{command}, \texttt{args}, \texttt{env}, \texttt{url}, \texttt{headers}.

\subsection{Scopes}

\begin{tabular}{lll}
\hline
\textbf{Scope} & \textbf{Куди записується} & \textbf{Коли юзати} \\
\hline
\texttt{local} (default) & \texttt{\textasciitilde/.claude.json} (current project) & Експерименти \\
\texttt{project} & \texttt{.mcp.json} у root & Команді через VCS \\
\texttt{user} & \texttt{\textasciitilde/.claude.json} (всі проєкти) & Особисті утиліти \\
\hline
\end{tabular}

\textbf{Precedence:} local $>$ project $>$ user $>$ plugin $>$ claude.ai connectors.

\subsection{Управління}

\begin{lstlisting}[language=bash]
claude mcp list                    # усі сервери
claude mcp get <name>              # деталі
claude mcp remove <name>           # видалити
claude mcp reset-project-choices   # скинути approval для .mcp.json

# усередині Claude Code:
/mcp                               # status, OAuth login
\end{lstlisting}

\section{Auth}

\begin{tabular}{lll}
\hline
\textbf{Транспорт} & \textbf{Метод} & \textbf{Як} \\
\hline
STDIO & Env-var & \texttt{--env API\_KEY=xxx} \\
HTTP & Bearer & \texttt{--header "Authorization: Bearer xxx"} \\
HTTP & API key & \texttt{--header "X-API-Key: xxx"} \\
HTTP & OAuth 2.0 & \texttt{/mcp} усередині Claude Code \\
\hline
\end{tabular}

MCP \textbf{рекомендує OAuth} для HTTP-серверів.

\section{Debug}

\subsection{MCP Inspector}

Інтерактивний UI, без install:

\begin{lstlisting}[language=bash]
# локальний сервер
npx @modelcontextprotocol/inspector \
  uv --directory /path/to/server run server.py

# npm-published
npx -y @modelcontextprotocol/inspector npx @modelcontextprotocol/server-filesystem /Users/me/Desktop

# PyPI-published
npx @modelcontextprotocol/inspector uvx mcp-server-git --repository ~/repo
\end{lstlisting}

\textbf{Вкладки:} Server connection (transport, args, env), Resources, Prompts, Tools, Notifications (stderr-логи).

\subsection{6-крокова діагностика}

\begin{enumerate}
\item \textbf{\texttt{claude mcp list}} показує сервер?
\item \textbf{\texttt{/mcp}} у сесії: connected?
\item \textbf{\texttt{claude mcp get <name>}} $\to$ запусти команду руками, дивись stderr
\item \textbf{stdout-факап?} \texttt{grep -n 'print(' server.py}
\item \textbf{Inspector} викликає tool вручну?
\item \textbf{Schema крива?} Type hints + docstring \texttt{Args:} секція
\end{enumerate}

\subsection{Reconnect i output limits}

\begin{itemize}
\item \textbf{HTTP/SSE:} автоматичний reconnect 5 спроб з exponential backoff
\item \textbf{STDIO:} НЕ reconnect (subprocess) --- перезапусти Claude Code
\item \textbf{Output warning} при 10K токенів. Підняти: \texttt{MAX\_MCP\_OUTPUT\_TOKENS=50000}
\item \textbf{Startup timeout:} \texttt{MCP\_TIMEOUT=10000 claude} (мс)
\end{itemize}

\section{Production-чек-ліст}

\begin{itemize}[leftmargin=2em]
\item[$\square$] \texttt{FastMCP("name")} kebab-case, унікальне
\item[$\square$] Type hints + docstring на кожному tool/resource/prompt
\item[$\square$] \texttt{Args:} секція у docstring (потрапляє у JSON Schema)
\item[$\square$] stderr-only logging для STDIO; жодного raw \texttt{print(...)}
\item[$\square$] Env-vars для секретів (не hard-code)
\item[$\square$] Output trimming для tool-ів з великим output
\item[$\square$] Тестовано в Inspector --- кожен tool, resource, prompt
\item[$\square$] README з прикладом \texttt{.mcp.json} і списком env-vars
\item[$\square$] Версіонування через git tags якщо публічно
\end{itemize}

\section{Команди для вправ}

\subsection{Setup}

\begin{lstlisting}[language=bash]
uv venv && source .venv/bin/activate
uv add "mcp[cli]"
\end{lstlisting}

\subsection{Inspector}

\begin{lstlisting}[language=bash]
npx @modelcontextprotocol/inspector uv run server.py
\end{lstlisting}

\subsection{Підключення до Claude Code (вправа 4)}

\begin{lstlisting}[language=bash]
claude mcp add --scope user --transport stdio my-clock \
  -- uv --directory $(pwd) run server.py

claude mcp list
claude
# усередині сесії:
/mcp
\end{lstlisting}

\section{Resources}

\begin{itemize}
\item \href{https://modelcontextprotocol.io/docs/concepts/architecture}{modelcontextprotocol.io --- Architecture}
\item \href{https://modelcontextprotocol.io/docs/develop/build-server}{modelcontextprotocol.io --- Build server (Python)}
\item \href{https://github.com/modelcontextprotocol/python-sdk}{github --- python-sdk}
\item \href{https://modelcontextprotocol.io/docs/tools/inspector}{modelcontextprotocol.io --- Inspector}
\item \href{https://code.claude.com/docs/en/mcp}{code.claude.com --- MCP integration}
\item Цей репо: \texttt{workshops/07-mcp-servers/exercises/} і \texttt{solutions/}
\end{itemize}

\end{document}
